\documentclass[11pt]{article}
\usepackage[margin=1in]{geometry}
\usepackage{amsmath,amssymb,amsthm}
\usepackage{enumitem}
\usepackage{xcolor}
\usepackage{tcolorbox}
\usepackage{hyperref}

\hypersetup{colorlinks=true, linkcolor=blue}

\newcommand{\Var}{\operatorname{Var}}
\newcommand{\E}{\mathbb{E}}

\title{HW4 Q4 Review --- Relevant Lecture Material\\[4pt]
\large From Slides09: Bernoulli Processes III --- Law of Large Numbers}
\author{ORF309/MAT380}
\date{}

\begin{document}
\maketitle

\noindent All content below is drawn from \textbf{ORF309\_MAT380\_S26\_Slides09} (pp.\ 47--56). These are the concepts needed for Q4 parts (a) and (b).

\section*{Q4(a): Law of Large Numbers}

\subsection*{Setup: The Sample Average}

Given an original random variable $X$ with $\E[X]$ defined, we repeat the experiment $n$ times to get i.i.d.\ copies $X_1, X_2, \ldots, X_n$. The \textbf{sample average} is:
\[
\frac{1}{n}\sum_{k=1}^{n} X_k
\]
This is itself a random variable for every finite $n$.

\subsection*{Expectation of the Sample Average (Slides09, p.\ 14)}

\[
\E\!\left[\frac{1}{n}\sum_{k=1}^{n} X_k\right] = \frac{1}{n}\sum_{k=1}^{n}\E[X_k] = \frac{1}{n}\cdot n\cdot \E[X] = \E[X]
\]
The expectation of the sample average equals $\E[X]$ for \emph{any} $n$ (no limit needed).

\subsection*{Theorem: Law of Large Numbers (Slides09, p.\ 16--18)}

\begin{tcolorbox}[colback=blue!5, colframe=blue!50!black, title=Law of Large Numbers for Random Variables]
With probability 1:
\[
\lim_{n\to\infty} \frac{1}{n}\sum_{k=1}^{n} X_k = \E[X]
\]
\end{tcolorbox}

\noindent Key observations from slides:
\begin{itemize}[nosep]
    \item For every finite $n$, the sample average $\frac{1}{n}\sum_{k=1}^n X_k$ is \textbf{random}.
    \item In the limit $n\to\infty$, the sample average is \textbf{not random} --- it equals $\E[X]$.
\end{itemize}

\bigskip
\noindent\textit{Think about how the Monte Carlo estimator $\hat{Z}_N = \frac{1}{N}\sum_{k=1}^N Z_k$ relates to this theorem.}

%--------------------------------------------------------------
\section*{Q4(b): Variance and the Proof of LLN}

\subsection*{Definition of Variance (Slides09, p.\ 24)}

\[
\Var(X) = \E\!\left[\big(X - \E[X]\big)^2\right]
\]
A random variable becomes ``less random'' when it is closer to its expectation (i.e., the squared deviations shrink).

\subsection*{Key Variance Properties (Slides09, p.\ 25)}

\begin{tcolorbox}[colback=green!5, colframe=green!50!black, title=Variance Properties]
\begin{enumerate}[nosep]
    \item $\Var(X) \geq 0$
    \item $\Var(X) = 0$ if and only if $X$ is not random
    \item $\Var(X) = \E[X^2] - \big(\E[X]\big)^2$
    \item If $X \perp Y$, then $\Var(X + Y) = \Var(X) + \Var(Y)$
    \item $\Var(\alpha X) = \alpha^2\,\Var(X)$, for all $\alpha \in \mathbb{R}$
\end{enumerate}
\end{tcolorbox}

\subsection*{Variance of the Sample Average (Slides09, pp.\ 27--31)}

The proof of LLN computes the variance of $S_n/n$ (the sample average). The slides work through this for indicator RVs, but the same logic applies to any i.i.d.\ random variables. The key computation pattern is:

\[
\Var\!\left(\frac{1}{n}\sum_{k=1}^n X_k\right)
\;\xrightarrow{\text{Prop.\ 5}}\;
\frac{1}{n^2}\,\Var\!\left(\sum_{k=1}^n X_k\right)
\;\xrightarrow{\text{Prop.\ 4}}\;
\frac{1}{n^2}\sum_{k=1}^n \Var(X_k)
\]

Since the $X_k$ are i.i.d.\ with $\Var(X_k) = \Var(X)$:

\begin{tcolorbox}[colback=orange!5, colframe=orange!50!black, title=Variance of the Sample Average]
\[
\Var\!\left(\frac{1}{n}\sum_{k=1}^n X_k\right) = \frac{1}{n^2}\cdot n\cdot \Var(X) = \frac{\Var(X)}{n}
\]
\end{tcolorbox}

\subsection*{Connection to Mean Square Error}

Recall from the slides (p.\ 27) that the proof of LLN uses:
\[
\E\!\left[\left(\frac{S_n}{n} - p\right)^{\!2}\right] = \Var\!\left(\frac{S_n}{n}\right)
\]
This works because $\E[S_n/n] = p$, so the squared deviation from the true mean \emph{is} the variance. The same logic connects to the MSE definition in Q4(b):
\[
\text{MSE}(N) = \E\!\Big[\big(\E[Z] - \hat{Z}_N\big)^2\Big]
\]

\bigskip
\noindent\textit{Use the expectation result (sample average is unbiased) and the variance result above to compute MSE$(N)$ in terms of $\Var(Z)$ and $N$.}

\end{document}
